% ============================================================================
\section{Multiprocessors and Instruction-Level Parallelism}
\label{sec:mp}

Flynn's classification, UMA and NUMA, superpipelining, superscalar and VLIW
processors, dependences and issue/completion order, multithreading, cache coherence
(snoopy and directory protocols, MSI, MESI) and synchronization (spin locks,
barriers).

\begin{remark}
The course slides for this part are not in this repository; this section follows
H\&P (Ch.~3 and~5) and the topic list in the lecturer's review slides. Check terms
against the lecture slides.
\end{remark}

% ----------------------------------------------------------------------------
\subsection{What to know}

\subsubsection{Classifying parallel machines}
\begin{center}
\small
\begin{tabular}{@{}lp{11.2cm}@{}}
\toprule
\multicolumn{2}{@{}l}{\textbf{Flynn's classification} (instruction streams $\times$ data streams)} \\
\midrule
SISD & one instruction stream, one data stream: a uniprocessor (may still be pipelined or superscalar). \\
SIMD & one instruction applied to many data: vector processors, multimedia extensions (SSE, AVX), GPUs. \\
MISD & several instruction streams on one data stream: \textbf{no commercial machine}; not used in practice. \\
MIMD & independent processors on their own data: multicores, multiprocessors, clusters. The most general and common. \\
\bottomrule
\end{tabular}
\end{center}
\begin{itemize}[leftmargin=*]
  \item \textbf{UMA} (uniform memory access), also \textbf{SMP} or centralized
  shared memory: all processors share one memory with the same access time; simple,
  but the bus/memory limits scaling (a few dozen cores).
  \item \textbf{NUMA} (non-uniform memory access), also distributed shared memory
  (DSM): memory is distributed among the nodes but still one address space; local
  accesses are faster than remote ones; scales further. Average access time
  $=f_{local}\,t_{local}+(1-f_{local})\,t_{remote}$.
  \item \textbf{Message passing} (clusters, multicomputers): private memories,
  explicit send/receive.
\end{itemize}

\subsubsection{Instruction-level parallelism (ILP)}
\begin{itemize}[leftmargin=*]
  \item \textbf{Superpipelining}: split stages into more, shorter stages, giving a
  higher clock rate. Costs: larger branch and load-use penalties (in cycles), more
  latch overhead.
  \item \textbf{Superscalar}: issue several instructions per cycle (ideal
  $\CPI=1/n$, $\text{IPC}=n$); hardware checks dependences dynamically, statically
  scheduled (in order) or dynamically scheduled (out of order); binary compatible.
  \item \textbf{VLIW}: the \emph{compiler} packs independent operations into one
  long instruction with a fixed slot per functional unit; no dynamic hazard
  detection, so simpler hardware. Drawbacks: \texttt{nop}s in unused slots increase
  \textbf{code size}; latencies are built into the code, so binaries are not
  compatible across generations; cache misses stall the whole bundle. (EPIC/IA-64 is
  a VLIW variant.)
\end{itemize}

\subsubsection{Dependences and issue order}
\begin{itemize}[leftmargin=*]
  \item \textbf{Data (true) dependence}: $j$ uses a result of $i$ (RAW hazard).
  \textbf{Name dependences}: \emph{anti-dependence} ($j$ writes a register $i$ reads:
  WAR) and \emph{output dependence} (both write it: WAW); no value flows, only the
  register name is shared. \textbf{Control dependence}: on a branch outcome.
  \item \textbf{Register renaming} (Tomasulo's reservation stations, or physical
  register files) gives each new value a fresh register and \textbf{removes WAR and
  WAW}; RAW must still be respected.
\end{itemize}
\begin{center}
\small
\begin{tabular}{@{}p{4.4cm}p{9.8cm}@{}}
\toprule
Policy & Hazards to check \\
\midrule
In-order issue, in-order completion & RAW only (like the 5-stage pipeline); a slow instruction blocks all behind it. \\
In-order issue, out-of-order completion & RAW and \textbf{WAW} (a later, shorter instruction can write first). Precise exceptions become hard. \\
Out-of-order issue, out-of-order completion & RAW, WAW and \textbf{WAR} (a later instruction may write before an earlier one reads). Needs an instruction window, renaming, and a reorder buffer for in-order \emph{commit} (precise exceptions, speculation). \\
\bottomrule
\end{tabular}
\end{center}
Scoreboarding issues out of order and stalls on WAR/WAW; Tomasulo's algorithm renames
registers to remove them; hardware speculation adds a reorder buffer.

\subsubsection{Multithreading}
Several threads share one core; each has its own PC and registers.
\begin{itemize}[leftmargin=*]
  \item \textbf{Fine-grained}: switch threads every cycle (round robin), skipping
  stalled threads; hides short and long stalls but slows each single thread.
  \item \textbf{Coarse-grained}: switch only on long stalls (L2 miss); little impact
  on one thread, but the pipeline must refill on a switch, so short stalls are not
  hidden.
  \item \textbf{Simultaneous multithreading (SMT)}: on a superscalar out-of-order
  core, issue instructions \emph{from several threads in the same cycle}, exploiting
  \textbf{both thread-level and instruction-level parallelism} (Intel
  Hyper-Threading).
\end{itemize}

\subsubsection{Cache coherence}
\begin{itemize}[leftmargin=*]
  \item With private caches, two processors can hold different values for one
  address. A system is \textbf{coherent} if a read returns the last value written
  (writes to one location are seen in the same order by all processors).
  Consistency is about the order of accesses to \emph{different} locations.
  \item \textbf{Snoopy protocols}: every cache watches (snoops) a shared bus;
  broadcast; good for small UMA machines. \textbf{Write-invalidate} (the usual):
  a writer first invalidates the other copies. \textbf{Write-update}: broadcast the new
  value.
  \item \textbf{Directory protocols}: a directory (at each block's home node) records
  the block's state and which caches share it (a bit vector); requests go point to
  point to the home, which sends invalidations only to the sharers. Scales to NUMA.
  Directory states: uncached, shared, exclusive/modified.
  \item Coherence misses (the fourth C): \emph{true sharing} (the word is really
  communicated) and \emph{false sharing} (different words in the same block).
\end{itemize}

\begin{center}
\small
\begin{tabular}{@{}p{2.6cm}p{11.6cm}@{}}
\toprule
State & Meaning \\
\midrule
M (Modified) & the only copy, dirty (memory is stale); must supply data and write back. \\
E (Exclusive) & MESI only: the only copy, clean (same as memory). \\
S (Shared) & clean, possibly in other caches. \\
I (Invalid) & not present / not valid. \\
\bottomrule
\end{tabular}
\end{center}
\begin{center}
\small
\begin{tabular}{@{}lll@{}}
\toprule
Event & MSI & MESI \\
\midrule
Read miss & BusRd $\to$ S & BusRd $\to$ E if no other copy, else S \\
Write miss & BusRdX $\to$ M & BusRdX $\to$ M \\
Write hit in S & BusUpgr (invalidate) $\to$ M & BusUpgr $\to$ M \\
Write hit in E & -- & \textbf{silent} E $\to$ M (no bus transaction) \\
Snoop BusRd, in M & flush (write back, supply data) $\to$ S & flush $\to$ S \\
Snoop BusRd, in E & -- & $\to$ S \\
Snoop BusRdX/BusUpgr & M: flush $\to$ I; S $\to$ I & M: flush $\to$ I; E, S $\to$ I \\
Evict & M: write back; S: silent & M: write back; E, S: silent \\
\bottomrule
\end{tabular}
\end{center}
The \textbf{E state} lets a processor that read a private line write it later without
a bus transaction: the common case for single-threaded code and private data.

\subsubsection{Synchronization}
\begin{itemize}[leftmargin=*]
  \item Built on an \textbf{atomic read-modify-write}: test-and-set, atomic exchange,
  fetch-and-increment, or the \texttt{ll}/\texttt{sc} pair (Section~\ref{sec:isa}).
  \item \textbf{Spin lock}: a processor loops until it acquires the lock. Spin on a
  \emph{cached copy} (plain loads, ``test-and-test-and-set'') and attempt the atomic
  operation only when the lock looks free; this keeps the spinning off the bus. With
  \texttt{ll}/\texttt{sc}:
\begin{lstlisting}
lockit: ll   $t0, 0($s0)       # read lock
        bne  $t0, $zero, lockit # spin while it is held (1)
        addi $t0, $zero, 1
        sc   $t0, 0($s0)       # try to set it to 1
        beq  $t0, $zero, lockit # sc failed: retry
        ...                     # critical section
        sw   $zero, 0($s0)     # release
\end{lstlisting}
  \item \textbf{Barrier}: no thread continues until all $n$ have arrived. Simple
  version: a lock-protected counter; the last arrival resets the counter and sets
  \texttt{release}; the others spin on \texttt{release}. Reusing it in a loop can
  deadlock (a fast thread re-enters before the slow ones have seen
  \texttt{release}); the fix is a \textbf{sense-reversing barrier}: each thread
  flips a private \texttt{local\_sense}, and the last arrival sets
  \texttt{release = local\_sense}.
\end{itemize}

% ----------------------------------------------------------------------------
\subsection{Common traps}
\begin{enumerate}[leftmargin=*]
  \item MISD is the class with no practical machines; SIMD is widely used (vectors,
  GPUs).
  \item Renaming removes WAR and WAW, never RAW.
  \item WAW appears with out-of-order \emph{completion}; WAR needs out-of-order
  \emph{issue} (or reads after later writes).
  \item VLIW: the compiler, not the hardware, finds the parallelism.
  \item SMT issues from several threads in the \emph{same} cycle; fine-grained
  switches every cycle but issues from one thread per cycle.
  \item MESI E $\to$ M needs no bus transaction; S $\to$ M does.
  \item A read snoop on an M line forces a write-back (flush).
\end{enumerate}

% ----------------------------------------------------------------------------
\subsection{Lecturer's review questions}

\begin{exercise}[Lecturer 1]
In the MESI protocol, what is the principal purpose of the E state?
(a) to avoid a second bus transaction when a single-threaded program writes a line
it previously read; (b) to identify the core responsible for responding to BusRd
requests from other cores; (c) to avoid the need for write-back when evicting an
exclusive line whose value is still correct in memory; (d) to help identify the
oldest line in a set when a victim must be chosen.
\end{exercise}
\begin{solution}
\textbf{(a)}. A read with no other sharers loads the line in E; a later write changes
E $\to$ M silently. With MSI the line would be in S, and the write would need a
BusUpgr. (c) is also true of S lines, so it is not what E adds.
\end{solution}

\begin{exercise}[Lecturer 2]
Register renaming helps to eliminate: (a) control hazards; (b) structural hazards;
(c) RAW and WAR data hazards; (d) WAR and WAW data hazards.
\end{exercise}
\begin{solution}
\textbf{(d)}: renaming removes name dependences. RAW is a true dependence; the
value must still be produced before it is used.
\end{solution}

\begin{exercise}[Lecturer 4]
The main feature of simultaneous multithreading is that it exploits both:
(a) thread-level and instruction-level parallelism; (b) coarse-grained multithreading
and speculation; (c) fine-grained multithreading and multilevel cache hierarchies;
(d) all of the above.
\end{exercise}
\begin{solution}
\textbf{(a)}: several threads (TLP) issue in the same cycle to a wide out-of-order
core (ILP).
\end{solution}

\begin{exercise}[Lecturer 5]
In Flynn's classification, which model is not much used in practice?
(a) SISD; (b) SIMD; (c) MIMD; (d) MISD.
\end{exercise}
\begin{solution}
\textbf{(d)} MISD.
\end{solution}

\begin{exercise}[Lecturer 6]
A VLIW machine: (a) relies on the compiler to schedule all instructions;
(b) dynamically detects hazards to reduce their delay penalty; (c) improves binary
compatibility between generations of the architecture; (d) improves code size.
\end{exercise}
\begin{solution}
\textbf{(a)}. (b) describes superscalar out-of-order hardware; (c) and (d) are VLIW's
weaknesses (latencies fixed in the code, \texttt{nop}-filled slots).
\end{solution}

% ----------------------------------------------------------------------------
\subsection{Exam-style questions}
Solutions are in Section~\ref{sec:mp-sol}.

\begin{exercise}[MSI vs.\ MESI trace]\label{q:mp-mesi}
Three processors with private write-back caches share a snoopy bus. Block $X$ is
initially in no cache. For each access give the state of $X$ in P1, P2, P3 and the bus
action, under MSI and under MESI:
\begin{center}
P1 reads $X$; P2 reads $X$; P1 writes $X$; P3 reads $X$; P2 writes $X$; P1 reads $X$.
\end{center}
Then P1 reads block $Y$ (cached nowhere) and writes it. Count the bus transactions for
$Y$ under each protocol.
\end{exercise}

\begin{exercise}[Dependences and issue policies]\label{q:mp-dep}
\leavevmode
\begin{lstlisting}
I1: div  r1, r2, r3     # long latency
I2: add  r4, r1, r5
I3: sub  r5, r6, r7
I4: mul  r4, r6, r8
I5: and  r9, r5, r4
\end{lstlisting}
\begin{enumerate}[label=(\alph*),leftmargin=*]
  \item List all true, anti and output dependences.
  \item With in-order issue and in-order completion, what happens while I1 executes?
  \item With in-order issue and out-of-order completion, which new hazard must be
  checked? Give the instruction pair.
  \item With out-of-order issue, which hazards appear if I3 or I4 runs before I2?
  \item Rename the registers so only true dependences remain.
\end{enumerate}
\end{exercise}

\begin{exercise}[VLIW scheduling]\label{q:mp-vliw}
A 2-slot VLIW has one memory slot and one ALU slot per instruction word. A load's
result can be used 2 cycles after it issues; an ALU result 1 cycle after. Registers
are read before they are written within one word. Schedule
\begin{lstlisting}
lw   r1, 0(r10)
lw   r2, 4(r10)
add  r3, r1, r2
sw   r3, 8(r10)
addi r10, r10, 12
\end{lstlisting}
in as few words as possible. How many words and \texttt{nop}s are there? What would a
superscalar processor do differently?
\end{exercise}

\begin{exercise}[Spin locks and barriers]\label{q:mp-sync}
\leavevmode
\begin{enumerate}[label=(\alph*),leftmargin=*]
  \item In the \texttt{ll}/\texttt{sc} spin lock above, why spin on the value with
  \texttt{bne} (line (1)) before trying \texttt{sc}? What would happen on the bus if
  every processor executed an atomic exchange in a tight loop?
  \item Explain why the simple counter barrier can fail when used twice in a loop,
  and how sense reversal fixes it.
\end{enumerate}
\end{exercise}

\begin{exercise}[UMA, NUMA and multithreading]\label{q:mp-misc}
\leavevmode
\begin{enumerate}[label=(\alph*),leftmargin=*]
  \item On a NUMA machine a local access takes 20~ns and a remote one 100~ns. Find
  the average access time if 80\% of accesses are local, and if only 50\% are.
  \item A core runs two threads; thread A stalls on long cache misses, thread B has
  many short 1-cycle dependence stalls. Which multithreading style hides each kind of
  stall?
  \item Why does superpipelining raise the branch penalty in cycles?
\end{enumerate}
\end{exercise}

% ----------------------------------------------------------------------------
\subsection{Solutions}
\label{sec:mp-sol}

\begin{noot}{Exercise~\ref{q:mp-mesi}}
States are listed as P1 P2 P3.
\begin{center}
\small
\begin{tabular}{@{}lllll@{}}
\toprule
Access & MSI states & MSI bus & MESI states & MESI bus \\
\midrule
P1 rd $X$ & S I I & BusRd & \textbf{E} I I & BusRd \\
P2 rd $X$ & S S I & BusRd & S S I & BusRd (P1 E $\to$ S) \\
P1 wr $X$ & M I I & BusUpgr & M I I & BusUpgr \\
P3 rd $X$ & S I S & BusRd, P1 flushes & S I S & BusRd, P1 flushes \\
P2 wr $X$ & I M I & BusRdX & I M I & BusRdX \\
P1 rd $X$ & S S I & BusRd, P2 flushes & S S I & BusRd, P2 flushes \\
\bottomrule
\end{tabular}
\end{center}
Both use 6 bus transactions for $X$ (with 2 write-backs). For $Y$: MSI uses BusRd
(S) then BusUpgr (M), 2 transactions; MESI uses BusRd (E) then a silent E $\to$ M,
1 transaction. This is the purpose of E.
\end{noot}

\begin{noot}{Exercise~\ref{q:mp-dep}}
\begin{enumerate}[label=(\alph*),leftmargin=*]
  \item True (RAW): I1$\to$I2 (r1), I3$\to$I5 (r5), I4$\to$I5 (r4). Anti (WAR):
  I2$\to$I3 (r5: I2 reads, I3 writes). Output (WAW): I2$\to$I4 (r4).
  \item I2 waits for r1, and because issue is in order, I3--I5 wait behind it even
  though I3 and I4 are independent of I1.
  \item WAW between I2 and I4: if I4 (\texttt{mul}) completes before I2, I2's later
  write would leave the wrong value in r4 for I5.
  \item I3 before I2: WAR, I3 overwrites r5 before I2 reads it. I4 before I2: WAW on
  r4.
  \item \texttt{I3: sub p5,r6,r7}; \texttt{I4: mul p4,r6,r8};
  \texttt{I5: and r9,p5,p4}. Now I3 and I4 can run at once during the \texttt{div};
  only I1$\to$I2 and I3, I4$\to$I5 remain.
\end{enumerate}
\end{noot}

\begin{noot}{Exercise~\ref{q:mp-vliw}}
\begin{center}
\small
\begin{tabular}{@{}lll@{}}
\toprule
Word & Memory slot & ALU slot \\
\midrule
1 & \texttt{lw r1,0(r10)} & \texttt{nop} \\
2 & \texttt{lw r2,4(r10)} & \texttt{nop} \\
3 & \texttt{nop} & \texttt{nop} \\
4 & \texttt{nop} & \texttt{add r3,r1,r2} \\
5 & \texttt{sw r3,8(r10)} & \texttt{addi r10,r10,12} \\
\bottomrule
\end{tabular}
\end{center}
5 words, 5 \texttt{nop}s (50\% of the slots). \texttt{r2} is ready in word 4, so
\texttt{add} cannot go earlier; \texttt{addi} shares word 5 with \texttt{sw} because
\texttt{sw} reads the old \texttt{r10}. A superscalar would find the same
parallelism at run time in hardware, with no \texttt{nop}s stored in the code, and the
same binary would still run on a wider or differently timed implementation.
\end{noot}

\begin{noot}{Exercise~\ref{q:mp-sync}}
\begin{enumerate}[label=(\alph*),leftmargin=*]
  \item While the lock is held, \texttt{ll}/\texttt{bne} just reads the cached copy
  (state S), so spinning causes no bus traffic. Only when the owner releases the lock
  (an invalidation) do the waiters miss, see 0 and try \texttt{sc}. A tight loop of
  atomic exchanges writes on every attempt, so the block ping-pongs between caches in
  M state with an invalidation on each try, flooding the bus and slowing the owner.
  \item The last thread resets the counter and sets \texttt{release=1}. A fast
  thread may leave and reach the barrier again, resetting \texttt{release=0} before a
  slow thread has seen the 1; the slow thread spins forever, and the fast threads
  wait for it: deadlock. With sense reversal each episode waits for the opposite
  value of \texttt{release}, so the next episode never undoes the previous release.
\end{enumerate}
\end{noot}

\begin{noot}{Exercise~\ref{q:mp-misc}}
\begin{enumerate}[label=(\alph*),leftmargin=*]
  \item $0.8\times20+0.2\times100=36$~ns; $0.5\times20+0.5\times100=60$~ns. Data
  placement matters on NUMA.
  \item Thread A's long misses: coarse-grained switching (or any style). Thread B's
  short stalls: only fine-grained or SMT, since coarse-grained switching costs a
  pipeline refill that is longer than the stall.
  \item The branch is resolved in a stage further from IF, so more stages (cycles)
  of wrongly fetched instructions must be flushed, even though each cycle is
  shorter.
\end{enumerate}
\end{noot}

